\section{Dữ liệu và chuẩn bị}

\subsection{Bốn nguồn dữ liệu}

Bốn nguồn ghi doanh thu ở bốn mức chi tiết khác nhau, nên không thể ghép từng dòng của nguồn này với từng dòng của nguồn kia (Bảng~\ref{tab:nguon}). Dữ liệu là bộ giả lập dựng từ Contoso Data Generator V2 \cite{contoso_v2}.

\captab{Bốn nguồn dữ liệu: một dòng là gì, khóa, ngày, tiền, dấu}\label{tab:nguon}\nopagebreak[4]
{\small
\begin{longtable}{@{}L{2.6cm} L{3.4cm} L{2.4cm} L{3.0cm} L{3.4cm}@{}}
\toprule
Nguồn & Một dòng là & Khóa & Ngày & Tiền, dấu \\
\midrule
\endhead
POS (Point of Sale, bán tại cửa hàng) & 1 món hàng trên 1 hóa đơn & \texttt{receipt\_no} + \texttt{line\_no} & Ngày bán & Nội tệ của nước có cửa hàng; dương \\
\midrule
WEB (đơn online) & 1 đơn hàng & \texttt{web\_order\_id} & Ngày đặt, ngày giao, ngày trả (kiểu Mỹ MM/DD/YYYY) & Tiền khách trả; dương \\
\midrule
Sổ cái & 1 bút toán, gộp cửa hàng $\times$ ngày ghi sổ hoặc ngày giao $\times$ loại tiền & \texttt{journal\_id} & Ngày ghi sổ & Nội tệ và USD theo tỷ giá cuối tháng; doanh thu âm, hàng trả dương \\
\midrule
Flash & 1 tháng $\times$ 1 kênh $\times$ 1 thị trường & Không có & Tháng đặt hàng & USD theo tỷ giá ngày đặt, làm tròn đến USD; dương \\
\bottomrule
\end{longtable}}

\textbf{Ai giữ nguồn nào.} File POS và file đơn online được xuất từ hệ thống bán hàng tại cửa hàng và hệ thống quản lý đơn online (OMS -- Order Management System); bộ phận vận hành bán hàng chịu trách nhiệm trạng thái đơn, ngày giao, hàng trả và việc xuất file. Sổ cái do phòng Tài chính giữ, gồm cả bút toán tay, bút toán bổ sung và quy đổi tỷ giá. Flash do phòng Kinh doanh lập ngày 05/01/2026; file đơn online được trích xuất ngày 15/01/2026. Khi một con số lệch, bảng này cho biết cần hỏi bộ phận nào.

\subsection{Mỗi bên ghi doanh thu thế nào: năm lý do lệch}

Hai bên đều đúng theo cách của mình nhưng khác nhau ở năm điểm (đề gọi là ``năm trục khác biệt''). Mục~IV chứng minh từng lý do bằng số trên dữ liệu.

\captab{Quy tắc ghi nhận của từng bên}\nopagebreak[4]
{\small
\begin{longtable}{@{}L{2.6cm} L{4.0cm} L{4.2cm} L{4.0cm}@{}}
\toprule
Lý do lệch & Flash & Sổ cái -- cửa hàng & Sổ cái -- online \\
\midrule
\endhead
Thời điểm & Ngày khách đặt hàng & Ngày ghi sổ; bán ngày nghỉ thì ghi ngày làm việc kế tiếp & Ngày giao thành công; hàng trả ghi giảm vào ngày nhận hàng trả \\
Giá trị & Giá niêm yết, trước giảm giá & Sau giảm giá & Sau giảm giá \\
Phạm vi & Gồm cả đơn sau đó bị hủy, không trừ hàng trả & Mọi hóa đơn POS; có thêm bút toán tay chuyển 3.480 USD của một cửa hàng sang online (tháng 6) & Không ghi đơn hủy, trừ hàng trả, có thêm bút toán tay phí vận chuyển \\
Tiền tệ & Tỷ giá ngày đặt & Tỷ giá cuối tháng ghi sổ & Tỷ giá cuối tháng ghi sổ \\
Mức chi tiết & Tháng $\times$ kênh $\times$ thị trường & Cửa hàng $\times$ ngày ghi sổ & Ngày giao $\times$ loại tiền \\
\bottomrule
\end{longtable}}

\subsection{Nạp dữ liệu bằng Power Query}

Toàn bộ dữ liệu được nạp và xử lý bằng Power Query (công cụ nạp và xử lý dữ liệu có sẵn trong Excel), theo luồng ở Hình~\ref{fig:luong}.

\begin{figure}[H]
\centering
\begin{tikzpicture}[node distance=0.55cm and 0.55cm, every node/.style={font=\small},
  hop/.style={draw=aivnGreen, rounded corners=3pt, fill=green!4, align=center, minimum height=1.1cm, text width=2.55cm},
  mui/.style={-{Stealth[length=2.2mm]}, thick, color=black!60}]
\node[hop] (tep) {4 file dữ liệu\\+ file danh mục\\(thư mục \texttt{data})};
\node[hop, right=of tep] (nap) {8 truy vấn nạp\\\texttt{q\_POS}, \texttt{q\_WEB},\\\texttt{q\_GL}, \texttt{q\_FLASH}…};
\node[hop, right=of nap] (gd) {Bảng giao dịch\\chung\\\texttt{q\_Transactions}};
\node[hop, right=of gd] (ds) {Đối soát\\cửa hàng, online,\\hàng trả};
\node[hop, below=of ds] (cn) {Bảng cầu nối\\\texttt{q\_BR\_Long},\\nhật ký ngoại lệ};
\node[hop, below=of nap] (kt) {Kiểm tra tổng,\\5 bẫy, chênh lệch\\tháng và kênh};
\draw[mui] (tep) -- (nap); \draw[mui] (nap) -- (gd); \draw[mui] (gd) -- (ds); \draw[mui] (ds) -- (cn);
\draw[mui] (nap) -- (kt);
\end{tikzpicture}

\capfig{Luồng xử lý trong file Power Query. Mỗi truy vấn là một file \texttt{.m} (file code của Power Query) trong thư mục \texttt{excel/queries}.}\label{fig:luong}
\end{figure}

Bốn quyết định khi nạp:
\begin{itemize}
  \item \textbf{Một ô đường dẫn thư mục.} Mọi truy vấn đọc ô \texttt{FolderPath} rồi tìm file theo phần đầu tên, không ghi cứng đường dẫn hay năm. Thư mục phải có đúng 1 file khớp tên; nếu không, truy vấn báo lỗi rõ ràng thay vì lấy bừa.
  \item \textbf{Đổi kiểu theo vùng en-US.} Mọi bước đổi kiểu số và ngày ghi rõ vùng English -- United States, để máy cài tiếng Việt vẫn đọc đúng.
  \item \textbf{Giữ dòng trùng, chỉ đánh dấu.} 12 dòng POS bị xuất trùng được giữ lại kèm cờ và lý do, để số dòng khớp bảng tổng kiểm soát và để đo được số tiền của dòng trùng.
  \item \textbf{Code truy vấn là bản gốc.} File Excel được dựng lại từ các file \texttt{.m} bằng một script, mọi bước có tên tiếng Việt và ghi chú.
\end{itemize}

\subsection{Năm bẫy khi nạp và cách xử lý}

Bảng~\ref{tab:bay} lấy thẳng từ sheet \texttt{Bay\_du\_lieu}: mọi số được tính lại mỗi lần làm mới dữ liệu.

\needspace{14\baselineskip}% dòng đầu của bảng cao khoảng 9 dòng chữ: giữ chú thích đi cùng
\captab{Năm bẫy: số trước và sau khi xử lý}\label{tab:bay}\nopagebreak[4]
{\small
\begin{longtable}{@{}p{0.4cm} p{2.6cm} p{3.4cm} p{4.2cm} p{4.2cm}@{}}
\toprule
\# & Bẫy & Truy vấn và bước xử lý & Trước khi xử lý & Sau khi xử lý \\
\midrule
\endhead
1 & Ngày kiểu Mỹ MM/DD/YYYY trong file đơn online & \texttt{q\_WEB}: Ô trống thành null; Đổi kiểu dữ liệu (ngày kiểu Mỹ, en-US) & Nếu đọc ngày đặt hàng theo kiểu ngày/tháng của Việt Nam: 5.950 đơn báo lỗi (phần ngày lớn hơn 12) và 3.215 đơn bị đảo ngày với tháng mà không báo lỗi & 9.441 đơn đọc đúng, 0 dòng lỗi; ngày đặt từ 01/12/2024 đến 31/12/2025 \\ \midrule
2 & File Flash có 3 dòng tiêu đề, 1 dòng trống và dòng TỔNG CỘNG & \texttt{q\_FLASH}: Bỏ các dòng tiêu đề; Bỏ dòng TỔNG CỘNG & 162 dòng thô; cộng cả cột doanh số ra 62.248.802 USD (gấp đôi vì có dòng TỔNG CỘNG) & 156 dòng chi tiết, tổng 31.124.401 USD \\ \midrule
3 & Sổ cái ghi doanh thu là số âm (ghi Có), hàng trả là số dương (ghi Nợ) & \texttt{q\_GL}: Đổi dấu doanh thu (USD) & Tổng amount\_usd = $-$28.698.496,23 USD & Tổng revenue\_usd = 28.698.496,23 USD \\ \midrule
4 & Dòng POS bị xuất trùng (trùng khóa receipt\_no + line\_no) & \texttt{q\_POS}: Sắp theo giờ xuất; Lần xuất đầu tiên của mỗi khóa; Cờ dòng trùng; Lý do trùng & 17.124 dòng, trong đó 12 dòng lặp khóa: cửa hàng POS 74, ngày bán 29/03/2025, xuất lại lúc 30/03/2025 08:42 & 17.112 dòng không trùng; dòng trùng có cờ is\_duplicate kèm lý do, không bị xóa \\ \midrule
5 & Dữ liệu có nhiều loại tiền & \texttt{q\_Transactions}: POS: ghép tỷ giá ngày bán; Đổi ra USD theo tỷ giá ngày (chia cho tỷ giá rồi mới cộng) & Cộng thẳng line\_amount của 5 loại tiền ra 14.706.296,35 (con số không có nghĩa) & Đổi từng dòng ra USD theo tỷ giá ngày bán rồi cộng: 14.060.501,00 USD (kể cả 12 dòng trùng, giống Flash); Flash kênh cửa hàng: 14.060.500 USD \\ \midrule
\end{longtable}}


Bẫy nguy hiểm nhất là ngày kiểu Mỹ: nếu đọc theo kiểu ngày/tháng của Việt Nam, phần lớn đơn báo lỗi, nhưng trong số còn lại có hơn ba nghìn đơn \textbf{đọc được mà sai tháng}, không có cảnh báo nào \nguon{QC}.

\subsection{Kiểm tra tổng}

Sau khi nạp, sheet \texttt{Kiem\_tra} tự so 13 chỉ tiêu với bảng tổng kiểm soát QC\_Summary của đề: \SoQCDat/\SoQCTong{} chỉ tiêu khớp, chênh lệch bằng 0 \nguon{QC}.

\captab{Kiểm tra tổng so với QC\_Summary}\nopagebreak[4]
{\small\begin{longtable}{@{}r p{1.1cm} p{5.0cm} >{\raggedleft\arraybackslash}p{2.5cm} >{\raggedleft\arraybackslash}p{2.5cm} r c@{}}
\toprule
STT & Nguồn & Chỉ tiêu & Đề (QC\_Summary) & Nhóm & Lệch & Kết quả \\
\midrule
\endhead
1 & POS & Số dòng & 17.124 & 17.124 & 0,00 & Đạt \\
2 & POS & Số hóa đơn khác nhau (receipt\_no) & 7.130 & 7.130 & 0,00 & Đạt \\
3 & POS & Khoảng ngày business\_date & 2025-01-01 $\rightarrow$ 2025-12-31 & 2025-01-01 $\rightarrow$ 2025-12-31 & 0,00 & Đạt \\
4 & POS & Tổng qty & 53.872 & 53.872 & 0,00 & Đạt \\
5 & POS & Tổng line\_amount (cộng thẳng mọi loại tiền, chỉ để kiểm tra nạp đủ) & 14.706.296,35 & 14.706.296,35 & 0,00 & Đạt \\
6 & POS & Tổng discount\_amount (cộng thẳng mọi loại tiền) & 872.455,82 & 872.455,82 & 0,00 & Đạt \\
7 & WEB & Số dòng (số đơn) & 9.441 & 9.441 & 0,00 & Đạt \\
8 & WEB & Số đơn theo trạng thái Delivered / Cancelled / Returned & 9121 / 139 / 181 & 9121 / 139 / 181 & 0,00 & Đạt \\
9 & WEB & Tổng net\_amount (cộng thẳng mọi loại tiền) & 18.712.884,87 & 18.712.884,87 & 0,00 & Đạt \\
10 & Sổ cái & Số bút toán & 6.940 & 6.940 & 0,00 & Đạt \\
11 & Sổ cái & Tổng amount\_usd (theo dấu kế toán) & $-$28.698.496,23 & $-$28.698.496,23 & 0,00 & Đạt \\
12 & Flash & Số dòng chi tiết (không tính tiêu đề và dòng tổng) & 156 & 156 & 0,00 & Đạt \\
13 & Flash & Tổng doanh số (USD) & 31.124.401 & 31.124.401 & 0,00 & Đạt \\
\bottomrule
\end{longtable}
}

\subsection{Bảng giao dịch chung và kiểm tra chất lượng dữ liệu}

Truy vấn \texttt{q\_Transactions} gộp \SoDongPOS{} dòng POS và \SoDongWEB{} đơn online thành một bảng \SoDongGiaoDich{} dòng cùng một dạng: mã cửa hàng bên kế toán (nối qua \texttt{store\_master}), ngày bán hoặc ngày đặt, ngày sổ cái ghi nhận (ngày ghi sổ hoặc ngày giao), ngày nhận hàng trả, tiền trước giảm giá, chiết khấu, sau giảm giá, và số USD theo cả hai tỷ giá: tỷ giá ngày (cách của Flash) và tỷ giá cuối tháng (cách của sổ cái). Mọi bảng đối soát và bảng cầu nối ở mục~IV chỉ đọc từ bảng này.

Mỗi lần làm mới, sheet \texttt{Chat\_luong\_du\_lieu} tự chạy \SoKiemTraChatLuong{} phép kiểm tra (Bảng~\ref{tab:chatluong}). \SoKiemTraCanXem{} phép ra kết quả ``Cần xem'', ứng với 3 lỗi đã biết và được ghi vào nhật ký ngoại lệ: dòng POS bị xuất trùng, một ngày bán của cửa hàng NX057 chưa có bút toán, và bút toán ADJ-0925 (vừa không gắn cửa hàng, vừa không có giấy tờ, nên vi phạm 2 phép) \nguon{QC}. Các phép còn lại bằng 0.

\captab{Kiểm tra chất lượng dữ liệu}\label{tab:chatluong}\nopagebreak[4]
{\small
\begin{longtable}{@{}r L{6.4cm} r c L{4.6cm}@{}}
\toprule
STT & Phép kiểm tra & Số lượng & Trạng thái & Chi tiết \\
\midrule
\endhead
1 & Dòng POS trùng khóa receipt\_no + line\_no & 12 & \textbf{Cần xem} & NX007 ngày 29/03/2025 \\
2 & Ô cửa hàng $\times$ ngày ghi sổ có bán mà sổ cái thiếu bút toán & 1 & \textbf{Cần xem} & NX057 ngày 21/11/2025 \\
3 & Bút toán sổ cái không gắn cửa hàng & 1 & \textbf{Cần xem} & JE20250904892 (MJ-ADJ-2509) \\
4 & Đơn hủy mà có ngày giao & 0 & Đạt &  \\
5 & Đơn giao trước ngày đặt & 0 & Đạt &  \\
6 & Hàng trả trước ngày giao & 0 & Đạt &  \\
7 & Dòng POS bán sau ngày cửa hàng đóng cửa & 0 & Đạt &  \\
8 & Dòng thiếu tỷ giá ngày (không đổi được ra USD) & 0 & Đạt &  \\
9 & Dòng POS có mã cửa hàng không có trong danh mục & 0 & Đạt &  \\
10 & Bút toán tay điều chỉnh không có giấy tờ (mã MJ-ADJ) & 1 & \textbf{Cần xem} & JE20250904892 30/09/2025 1.250,00 USD \\
11 & Dòng thiếu tỷ giá cuối tháng (tháng ghi nhận hoặc tháng nhận hàng trả) & 0 & Đạt &  \\
12 & Dòng POS không tìm được ngày ghi sổ (ngày bán ngoài lịch) & 0 & Đạt &  \\
13 & Số dòng bảng giao dịch chung $-$ (dòng POS + đơn online) & 0 & Đạt &  \\
\bottomrule
\end{longtable}}


\subsection{Cách tính thứ hai: mô hình Power BI}

Để kiểm tra lại số của Power Query, nhóm dựng mô hình Power BI trên file dữ liệu mentor đã nối sẵn (\texttt{Nexa\_Home\_Recon\_Dataset.xlsx}): 12 bảng, gồm 6 bảng danh mục (có bảng ngày), 4 bảng giao dịch và 2 bảng tỷ giá. Ngày ghi nhận của sổ cái (ngày ghi sổ, ngày giao) là quan hệ chủ động; ngày bán, ngày đặt, ngày trả là quan hệ bị động, chỉ bật bằng hàm \texttt{USERELATIONSHIP} khi cần. Số tổng tính bằng DAX (Data Analysis Expressions, ngôn ngữ công thức của Power BI) khớp với Power Query ở cả năm, 12 tháng và 3 kênh. Phần so từng khoản cầu nối ở mục~IV.5.
